% Options for packages loaded elsewhere
\PassOptionsToPackage{unicode}{hyperref}
\PassOptionsToPackage{hyphens}{url}
%
\documentclass[
]{article}
\usepackage{amsmath,amssymb}
\usepackage{iftex}
\ifPDFTeX
  \usepackage[T1]{fontenc}
  \usepackage[utf8]{inputenc}
  \usepackage{textcomp} % provide euro and other symbols
\else % if luatex or xetex
  \usepackage{unicode-math} % this also loads fontspec
  \defaultfontfeatures{Scale=MatchLowercase}
  \defaultfontfeatures[\rmfamily]{Ligatures=TeX,Scale=1}
\fi
\usepackage{lmodern}
\ifPDFTeX\else
  % xetex/luatex font selection
\fi
% Use upquote if available, for straight quotes in verbatim environments
\IfFileExists{upquote.sty}{\usepackage{upquote}}{}
\IfFileExists{microtype.sty}{% use microtype if available
  \usepackage[]{microtype}
  \UseMicrotypeSet[protrusion]{basicmath} % disable protrusion for tt fonts
}{}
\makeatletter
\@ifundefined{KOMAClassName}{% if non-KOMA class
  \IfFileExists{parskip.sty}{%
    \usepackage{parskip}
  }{% else
    \setlength{\parindent}{0pt}
    \setlength{\parskip}{6pt plus 2pt minus 1pt}}
}{% if KOMA class
  \KOMAoptions{parskip=half}}
\makeatother
\usepackage{xcolor}
\usepackage[margin=1in]{geometry}
\usepackage{graphicx}
\makeatletter
\newsavebox\pandoc@box
\newcommand*\pandocbounded[1]{% scales image to fit in text height/width
  \sbox\pandoc@box{#1}%
  \Gscale@div\@tempa{\textheight}{\dimexpr\ht\pandoc@box+\dp\pandoc@box\relax}%
  \Gscale@div\@tempb{\linewidth}{\wd\pandoc@box}%
  \ifdim\@tempb\p@<\@tempa\p@\let\@tempa\@tempb\fi% select the smaller of both
  \ifdim\@tempa\p@<\p@\scalebox{\@tempa}{\usebox\pandoc@box}%
  \else\usebox{\pandoc@box}%
  \fi%
}
% Set default figure placement to htbp
\def\fps@figure{htbp}
\makeatother
\setlength{\emergencystretch}{3em} % prevent overfull lines
\providecommand{\tightlist}{%
  \setlength{\itemsep}{0pt}\setlength{\parskip}{0pt}}
\setcounter{secnumdepth}{-\maxdimen} % remove section numbering
\usepackage{bookmark}
\IfFileExists{xurl.sty}{\usepackage{xurl}}{} % add URL line breaks if available
\urlstyle{same}
\hypersetup{
  pdftitle={Support Vector Regression for Coffee Quality Prediction},
  pdfauthor={Luisa Kalkert},
  hidelinks,
  pdfcreator={LaTeX via pandoc}}

\title{Support Vector Regression for Coffee Quality Prediction}
\author{Luisa Kalkert}
\date{2025-12-28}

\begin{document}
\maketitle

\subsection{Support Vector Regression}\label{support-vector-regression}

\subsubsection{Data Preprocessing and
Cleaning}\label{data-preprocessing-and-cleaning}

\paragraph{Feature selection:}\label{feature-selection}

In SVM Regression, categorical variables are One-Hot Encoded using dummy
variables. This means, features with many unique values will lead to a
sparse feature matrix. This is disadvantageous as it increases
computational cost and poses the risk of overfitting. Additionally,
categories with only a few or even only one observation most likely will
not add to model robustness, but are adding noise to the signal.
Therefore, we are removing categorical features with many unique values,
only keeping columns with a low number of unique values. E.g. we are not
including the ``Region'' column with 344 unique values in the model
training, because we'd on average have less than 4 observations per
Region. For features with a high number of missing values, we either
remove the column or impute the missing values. For instance for the
``Variety'' column, we have about 15\% missing values and 8\% entries
``other'' at 30 different varieties. Keeping those, would mean adding 30
sparse columns, and at the same time the information gain is most likely
not too high with this many missing values. Therefore, we are not
including this feature.

\paragraph{Imputation:}\label{imputation}

For the altitude column, we have about 17\% missing values. Here, we are
imputing the missing values from the ``Region'' column. For each
observation with missing altitude, we are imputing the median altitude
of all coffees from the same region. This covers almost all missing
values. For the remaining observations, we are imputing the mean
altitude of all coffees from the same country of origin. As the
``Region'' and ``Country'' columns are not used in the model training,
we are not creating correlated features with this imputation. We are
imputing the missing values after the train/test split, using only
values from the training set, to avoid data leakage from the test set
into the training set.

\paragraph{Further preprocessing:}\label{further-preprocessing}

The bag weight column contains both pounds and kilograms, so we are
converting all to kilograms. However, for some entries, it is unclear
which unit is used. For these, we are averaging the weight in pounds and
kilograms. This means for those entries the weight is certainly
incorrect, but should be within an order of magnitude of the correct
weight. If we can assume that speciality coffee is sold in smaller
quantities, we can still gain valuable information from the weight.

\paragraph{Scaling:}\label{scaling}

The SVM method used for model training has a built in scaling
functionality. Both predictor and target variables are scaled to have
expected value 0 and variance 1.
{[}\url{https://rdrr.io/rforge/e1071/man/svm.html}{]} This
transformation is based on the training data and applied to train and
test data sets.

\subsubsection{Support Vector Regression - Hyperparameter
Optimization}\label{support-vector-regression---hyperparameter-optimization}

For the prediction of cupper points we train an SVM for regression on
the data set, holding out the test data set for final evaluation. To
choose appropriate hyperparameters, we conduct a hyperparameter
optimization on the training data set using 10-fold cross-validation. We
use a grid search for different values for cost and epsilon for a
radial, linear, and 2-degree polynomial kernel. We compare the mean
absolute error between runs to choose the best-performing hyperparameter
combination. Finally, we train a model with these hyperparameters on the
entire training set and evaluate on the test data set.

\pandocbounded{\includegraphics[keepaspectratio]{arabica_SVM_files/figure-latex/unnamed-chunk-3-1.pdf}}
These figures show the mean absolute error for the different
hyperparameters by kernel. The best performing hyperparameter
combination for each kernel are marked with orange points.

\begin{verbatim}
## Best SVM kernel and parameters selected:
\end{verbatim}

\begin{verbatim}
##   kernel epsilon      cost
## 2 linear     0.1 0.0078125
\end{verbatim}

\begin{verbatim}
## [1] "MAE: 0.143099275160558"
\end{verbatim}

The best performing kernel is the linear kernel, which means the
relationship between the features and the target variable is best
approximated by a multidimensional hyperplane.

\pandocbounded{\includegraphics[keepaspectratio]{arabica_SVM_files/figure-latex/unnamed-chunk-5-1.pdf}}

Here, we can see the predictions vs.~the true values for the test data
set. We can see that aside from two outliers, the predictions are close
to the true values.

\end{document}
